\documentclass[11pt]{article}
\usepackage[localtheorems]{raeez-math-template}

%%% Packages

%%% Page geometry

%%% Theorem environments
\theoremstyle{plain}
\newtheorem{theorem}{Theorem}[section]
\newtheorem{proposition}[theorem]{Proposition}
\newtheorem{corollary}[theorem]{Corollary}
\newtheorem{lemma}[theorem]{Lemma}
\newtheorem{conjecture}[theorem]{Conjecture}

\theoremstyle{definition}
\newtheorem{definition}[theorem]{Definition}
\newtheorem{construction}[theorem]{Construction}
\newtheorem{example}[theorem]{Example}

\theoremstyle{remark}
\newtheorem{remark}[theorem]{Remark}

%%% Macros
\newcommand{\fg}{\mathfrak{g}}
\newcommand{\fh}{\mathfrak{h}}
\newcommand{\fgl}{\mathfrak{gl}}
\newcommand{\fsl}{\mathfrak{sl}}
\newcommand{\fL}{\mathfrak{L}}
\newcommand{\Ran}{\mathrm{Ran}}
\newcommand{\Mbar}{\overline{\mathcal{M}}}
\newcommand{\FM}{\mathrm{FM}}
\newcommand{\MC}{\mathrm{MC}}
\newcommand{\cA}{\mathcal{A}}
\newcommand{\cC}{\mathcal{C}}
\newcommand{\cH}{\mathcal{H}}
\newcommand{\cM}{\mathcal{M}}
\newcommand{\cO}{\mathcal{O}}
\newcommand{\cW}{\mathcal{W}}
\newcommand{\cZ}{\mathcal{Z}}
\newcommand{\bC}{\mathbb{C}}
\newcommand{\bP}{\mathbb{P}}
\newcommand{\bR}{\mathbb{R}}
\newcommand{\bS}{\mathbb{S}}
\newcommand{\bZ}{\mathbb{Z}}
\newcommand{\Ainf}{A_\infty}
\newcommand{\Eone}{E_1}
\newcommand{\Etwo}{E_2}
\newcommand{\En}{E_n}
\newcommand{\Sd}{\bS^d}
\newcommand{\Res}{\operatorname{Res}}
\newcommand{\Tr}{\operatorname{Tr}}
\newcommand{\HH}{\mathrm{HH}}
\newcommand{\HC}{\mathrm{HC}}
\newcommand{\CC}{\mathrm{CC}}
\newcommand{\Zder}{Z^{\mathrm{der}}_{\mathrm{ch}}}
\newcommand{\ChirHoch}{\mathrm{ChirHoch}}
\newcommand{\CoHA}{\mathrm{CoHA}}
\newcommand{\Fuk}{\mathrm{Fuk}}
\newcommand{\Rep}{\mathrm{Rep}}
\newcommand{\Ext}{\mathrm{Ext}}
\newcommand{\Fact}{\mathrm{Fact}}
\newcommand{\CY}{\mathrm{CY}}
\newcommand{\DT}{\mathrm{DT}}
\newcommand{\GW}{\mathrm{GW}}
\newcommand{\Hom}{\mathrm{Hom}}
\newcommand{\Conf}{\mathrm{Conf}}
\newcommand{\MF}{\mathrm{MF}}
\newcommand{\Coh}{\mathrm{Coh}}
\newcommand{\SCchtop}{\mathrm{SC}^{\mathrm{ch},\mathrm{top}}}
\newcommand{\SpCh}{\mathrm{Sp}^{\mathrm{ch}}}
\newcommand{\PhiFA}{\Phi^{\mathrm{FA}}}
\DeclareMathOperator{\Aut}{Aut}
\DeclareMathOperator{\Spec}{Spec}

%%% Title
\title{The Calabi--Yau-to-chiral functor\\
and the $\kappa_\bullet$-spectrum}
\author{Raeez Lorgat}
\date{}

\begin{document}
\maketitle

%%% ================================================================
%%% ABSTRACT
%%% ================================================================

\begin{abstract}
Where do chiral algebras come from?  Volumes~I and~II of the
modular Koszul duality programme accept a chiral algebra as given
and extract its invariants: bar-cobar adjunction, the modular
characteristic~$\kappa_{\mathrm{ch}}$, the shadow tower, the five
theorems A--D+H.  The geometric source is left unspecified.  This
paper constructs it.

A Calabi--Yau category $\cC$ of dimension~$d$ carries a
negative-cyclic Calabi--Yau class
$[\sigma_\cC]\in HC^{-}_{d}(\cC)$ whose Hochschild shadow is a
non-degenerate cyclic pairing.  On the verified specialisation locus,
the CY-to-chiral assignment
\[
 \Phi_d \;=\; \SpCh_{\Sigma_{d-1}, C} \circ \Phi^{\mathrm{FA}}_d
 \colon \CY_d^{\mathrm{wit}}\text{-}\mathrm{Cat} \longrightarrow
 \En\text{-}\mathrm{ChirAlg}
\]
factors in two stages. Stage~1, $\Phi^{\mathrm{FA}}_d$, produces a
holomorphic $E_d$-factorisation algebra
$\Phi^{\mathrm{FA}}_d(\cC)$ on a CY$_d$ variety~$X$ after the needed
formality, locality, and completion witnesses are fixed.  Its algebraic
input is the Gerstenhaber bracket of degree $-1$ on
$\HH^\bullet(\cC)$, or degree $1-d$ after CY duality to Hochschild
homology. Stage~2,
$\SpCh_{\Sigma_{d-1}, C}$, is factorisation homology over a
closed $(d-1)$-cycle $\Sigma_{d-1} \subset X$ followed by restriction
to a reference curve $C \subset X$. The four structural operations,
Lie conformal algebra $\fL_\cC$ from the Gerstenhaber bracket,
factorisation envelope $\Fact_C(\fL_\cC)$, witnessed framing data, and
negative-cyclic quantisation, realise the composite. The assignment is
proved at $d = 2$ (Theorem~CY-A$_2$, $\Etwo$-chiral on any curve~$C$).
At $d = 3$ it is an object-level $\Eone$-chiral construction on the
witnessed filtered locus. A single CY$_d$ category admits a family of
$\Eone$-chiral shadows indexed by the choice of $(\Sigma_{d-1}, C)$;
the Borcherds--$\Delta_5$, Igusa--$\chi_{10}$, and paramodular
Gritsenko--Cl\'ery outputs at $d = 3$ are distinct specialisations of
one canonical $\Phi^{\mathrm{FA}}_3$.

A single CY manifold admits multiple chiral algebraizations, each
with its own modular characteristic.  The resulting
$\kappa_\bullet$-spectrum
$\Spec_{\kappa_\bullet}(X) = \{\kappa_{\mathrm{cat}},
\kappa_{\mathrm{ch}}^{\mathrm{Heis}}, \kappa_{\mathrm{BKM}},
\kappa_{\mathrm{fiber}}\}$: four numbers measuring four faces of
the same manifold.  For $K3 \times E$:
$\Spec_{\kappa_\bullet} = \{0, 3, 5, 24\}$. The value~$2$ belongs to
the K3 fibre lane, not to the total-space categorical characteristic.
The shadow--Siegel
gap theorem identifies four structural obstructions separating
$\kappa_{\mathrm{ch}}^{\mathrm{Heis}} = 3$ from $\kappa_{\mathrm{BKM}} = 5$.

For toric CY$_3$, the functor interfaces with the critical
cohomological Hall algebra: $\CoHA(Q_X, W_X) \simeq
Y^+(\widehat{\fgl}_{Q_X})$ provides the $\Eone$-sector, and the
Drinfeld center recovers $\Etwo$-braiding.  The $\Eone$
stabilization theorem proves that for toric $GL(d)$-invariant local
models at $d \geq 3$, the relevant invariant Schouten--Nijenhuis
sector is silent. General CY categories acquire their $\Eone$ output
from Stage~2 and the CY degree shift.
\end{abstract}

\tableofcontents

%%% ================================================================
%%% 1. INTRODUCTION
%%% ================================================================

\section{Introduction}\label{sec:intro}

\subsection{The source problem}

A chiral algebra is a solution.  What is the problem it solves?

The bar-cobar machine of Volume~I extracts from any chiral
algebra~$A$ on a curve~$X$ a sequence of invariants: the bar
complex $B(A) = T^c(s^{-1}\bar{A})$ with its deconcatenation
coproduct, the universal Maurer--Cartan element $\Theta_A$
satisfying $d\Theta + \frac{1}{2}[\Theta, \Theta] = 0$, the
modular characteristic $\kappa_{\mathrm{ch}}(A)$ governing the
genus-$1$ curvature, and the shadow tower
$\{S_r\}_{r \geq 2}$ encoding higher-genus data.  The five
theorems (A--D+H) are structural properties of this machine.
Volume~II interprets the output physically: the derived chiral
center $\Zder(A) = C^\bullet_{\mathrm{ch}}(A,A)$ is the bulk of
a 3d holomorphic-topological gauge theory; the
$\SCchtop$ pair $(C^\bullet_{\mathrm{ch}}(A,A), A)$
is the boundary--bulk system.

Neither volume constructs a single chiral algebra.  Affine
Kac--Moody algebras arise from flat connections; Virasoro from
reparametrizations; every other family in the standard landscape
must be supplied by hand.  The question is: what functor connects
geometric input to chiral algebraic output, and what structure
must it preserve?

The answer is the CY-to-chiral functor
$\Phi_d = \SpCh_{\Sigma_{d-1}, C} \circ \Phi^{\mathrm{FA}}_d$.
A Calabi--Yau category $\cC$ carries a cyclic $\Ainf$-structure
encoding a $d$-dimensional Frobenius condition at the chain level.
The first stage $\Phi^{\mathrm{FA}}_d$ converts this datum into a
canonical $E_d$-holomorphic factorisation algebra on the CY$_d$
variety~$X$; the second stage $\SpCh_{\Sigma_{d-1}, C}$
specialises along a closed $(d-1)$-cycle $\Sigma_{d-1} \subset X$
and restricts to a reference curve $C \subset X$, producing the
chiral algebra $A_\cC$ on~$C$ with three preservation properties:
the bar complex of $A_\cC$ recovers the cyclic bar complex
of~$\cC$, the modular characteristic $\kappa_{\mathrm{ch}}(A_\cC)$
equals the CY Euler characteristic $\chi^{\CY}(\cC)$, and the
$\Sd$-framing determines the operadic level on~$C$.

\subsection{The $\kappa_\bullet$-spectrum}

A single CY manifold admits multiple chiral algebraizations.  The
lattice vertex algebra $V_\Lambda$ attached to $H^*(K3, \bZ)$,
the generalized Kac--Moody algebra whose denominator is the Igusa
cusp form $\Delta_5$, and the Conway module from the Leech lattice
frame are three algebraizations of the same K3 surface.  Each
carries its own modular characteristic.

The $\kappa_\bullet$-spectrum records all four lanes:
\begin{align*}
 \kappa_{\mathrm{cat}} &= \chi(\cO_S)
   &&\text{(holomorphic Euler characteristic)}, \\
 \kappa_{\mathrm{ch}}^{\mathrm{Heis}} &= \dim_\bC X
   &&\text{(Heisenberg--Mukai chiral specialisation)}, \\
 \kappa_{\mathrm{BKM}} &= \tfrac{1}{2}\operatorname{wt}(\Phi_{10}^{\mathrm{un}})
   &&\text{(Borcherds--Kac--Moody automorphic weight)}, \\
 \kappa_{\mathrm{fiber}} &= \operatorname{rank}(\Lambda)
   &&\text{(lattice rank)}.
\end{align*}
For $K3 \times E$:
$\Spec_{\kappa_\bullet} = \{0, 3, 5, 24\}$.  The four numbers
measure four faces of the same manifold: holomorphic Euler
characteristic, boundary chiral algebra, bulk BKM superalgebra,
and lattice frame.  The chiral-holographic pair
$(\kappa_{\mathrm{ch}}^{\mathrm{Heis}}, \kappa_{\mathrm{BKM}}) = (3, 5)$ governs
the boundary--bulk axis of Volumes~I--II.  The
$\kappa_\bullet$-spectrum is invisible from any single chiral
algebra; it is intrinsic to the CY category.

\subsection{The $\Eone$ stabilization theorem}

For $d \geq 3$, the CY chiral algebra is natively $\Eone$, not
$E_d$.  The full Schouten--Nijenhuis bracket on
$\mathrm{PV}^*(\bC^d)$ does not vanish.  The toric
$GL(d)$-invariant CY sector used below is the abelian local sector,
and the CoHA extension correspondence (sub before quotient) imposes
a preferred ordering, which is $\Eone$.  The braided monoidal structure is
recovered through the Drinfeld center of the $\Eone$-monoidal
representation category.

The obstruction to additional structure beyond $\Eone$ is
controlled by $\pi_d(BU)$, which is $\bZ$ for even~$d$ and
$0$ for odd~$d$ by Bott periodicity.  The effective obstruction is
trivial precisely at $d \bmod 8 \in \{1, 3, 7\}$.

\subsection{Outline}

Section~\ref{sec:cy-cats} reviews CY categories and cyclic
$\Ainf$-structures.  Section~\ref{sec:functor} constructs the
four-step functor $\Phi$ and proves Theorem~CY-A$_2$.
Section~\ref{sec:kappa-spectrum} defines the
$\kappa_\bullet$-spectrum and develops its properties.
Section~\ref{sec:k3e} works through the $K3 \times E$ prototype
in detail: the thirteen structural results K3-1 through K3-13,
the shadow--Siegel gap theorem, and the four obstructions
separating $\kappa_{\mathrm{ch}}$ from $\kappa_{\mathrm{BKM}}$.
Sections~\ref{sec:toric}--\ref{sec:e1-stabilization} treat toric
CY$_3$, the CoHA, and the $\Eone$ stabilization theorem.


%%% ================================================================
%%% 2. CY CATEGORIES AND CYCLIC A-INFINITY
%%% ================================================================

\section{Calabi--Yau categories and cyclic $\Ainf$-structures}
\label{sec:cy-cats}

\subsection{CY categories}

\begin{definition}[Calabi--Yau category]
\label{def:v3-st-cy-category}
A smooth proper dg category $\cC$ over a field~$k$ is
\emph{Calabi--Yau of dimension~$d$} if there exists a
quasi-isomorphism of $\cC$-bimodules
\[
 \cC^\vee \;\xrightarrow{\;\sim\;}\; \cC[d]
\]
where $\cC^\vee = \mathrm{RHom}_{\cC^e}(\cC, \cC^e)$ is the
inverse dualizing bimodule and $[d]$ denotes the cohomological
shift by~$d$.
\end{definition}

The CY condition encodes Serre duality at the categorical level.
For the bounded derived category $D^b(\Coh(X))$ of a smooth
projective variety~$X$ with trivial canonical bundle
$\omega_X \simeq \cO_X$, Serre duality gives
$\bS_X \simeq [d]$ on $D^b(\Coh(X))$ (where $d = \dim_\bC X$),
and the resulting $\CY_d$-structure is the classical Serre
functor.

\begin{example}[Standard families]
\label{ex:v3-st-cy-families}
Four families of CY categories provide the standard landscape:
\begin{enumerate}[(i)]
 \item \emph{Elliptic curves} ($d = 1$): $D^b(\Coh(E_\tau))$ is
 $\CY_1$; the functor $\Phi$ produces the Heisenberg vertex
 algebra $H_k$ at level $k = \mathrm{vol}(E_\tau)$.
 \item \emph{K3 surfaces} ($d = 2$): $D^b(\Coh(S))$ is $\CY_2$;
 $\Phi$ produces an $\Etwo$-chiral algebra with
 $\kappa_{\mathrm{ch}} = \chi(\cO_S) = 2$.
 \item \emph{CY threefolds} ($d = 3$): the conifold, local $\bP^2$,
 $K3 \times E$.  The functor $\Phi$ at $d = 3$ is proved
 (Theorem~CY-A$_3$).
 \item \emph{CY fourfolds} ($d = 4$): $K3 \times K3$, the sextic
 fourfold.  The obstruction $\pi_4(BU) = \bZ$ controls the
 shifted symplectic structure.
\end{enumerate}
\end{example}

\subsection{Cyclic $\Ainf$-structure}

A CY category $\cC$ of dimension~$d$ carries a negative-cyclic
Calabi--Yau class whose Hochschild shadow is a non-degenerate trace:
\[
 \Tr \colon \HH_d(\cC) \longrightarrow k.
\]
The negative-cyclic class is the CY datum; the Hochschild trace is
its decategorified shadow.  The primary chain object is the cyclic
bar complex $\CC_\bullet(\cC)$ equipped with the Connes
$B$-operator.  Higher framing data are additional witnesses, not a
formal consequence of the trace alone.

\begin{definition}[Cyclic $\Ainf$-structure]
\label{def:v3-st-cyclic-ainf}
A \emph{cyclic $\Ainf$-structure} on a CY$_d$ category $\cC$
consists of:
\begin{enumerate}[(i)]
 \item An $\Ainf$-category structure on $\cC$: operations
 $\mu_k \colon \cC^{\otimes k} \to \cC[2-k]$ for $k \geq 1$
 satisfying the Stasheff identities.
 \item A non-degenerate cyclic pairing
 $\langle -, - \rangle \colon \cC \otimes \cC \to k[-d]$
 satisfying the higher cyclicity conditions
 $\langle \mu_k(a_1, \ldots, a_k), a_0 \rangle
 = \pm \langle \mu_k(a_0, a_1, \ldots, a_{k-1}), a_k \rangle$.
 \item The induced Connes $B$-operator $B \colon \CC_n(\cC) \to
 \CC_{n+1}(\cC)$ with $B^2 = 0$ and $Bb + bB = 0$.
\end{enumerate}
\end{definition}

\begin{proposition}[Categorical Hochschild homology and the Gerstenhaber bracket]
\label{prop:v3-st-gerstenhaber}
The categorical Hochschild cohomology $\HH^\bullet(\cC)$ of any dg category
carries a Gerstenhaber bracket $[-,-]$ of degree~$-1$.  For a
$\CY_d$ category, the CY pairing identifies
$\HH^\bullet(\cC) \simeq \HH_{d - \bullet}(\cC)$, and the
Gerstenhaber bracket becomes a Lie bracket on $\HH_\bullet(\cC)$
of degree $1 - d$.
\end{proposition}

\begin{proof}
The Gerstenhaber bracket is a standard structure on categorical Hochschild
cohomology (Gerstenhaber 1963, Keller 2006); the CY identification
with homology uses the non-degenerate trace and the
Hochschild--Kostant--Rosenberg degeneration at the $E_1$-page of
the Hodge-to-de~Rham spectral sequence.
\end{proof}

\subsection{The framing witness}

The CY$_d$-structure on $\cC$ supplies the negative-cyclic and
Connes data.  A chain-level $\Sd$-framing is an additional witness.
The $\bS^1$-action on $\HH_\bullet(\cC)$ is the Connes shadow; the
full framing, when constructed, governs the operadic level of the
chiral algebra.

\begin{proposition}[Framing witness; Kontsevich--Vlassopoulos]
\label{prop:v3-st-sd-framing}
For a $\CY_d$ category $\cC$, the negative-cyclic CY class supplies
the Connes action on $\CC_\bullet(\cC)$.  At $d = 2$, the
Kontsevich--Vlassopoulos framing witness provides the $\bS^2$-action
needed for the $\Etwo$-algebra structure on cyclic homology.  At
$d = 3$, the corresponding chain-level $\bS^3$-framing is part of
the witnessed filtered locus.
\end{proposition}

The framing witness is the datum that the assignment $\Phi$ consumes
to determine the operadic enhancement.  At $d = 2$, the
$\bS^2$-framing is proved.  At $d = 3$, the chain-level
$\bS^3$-framing compatible with the BV structure is one of the
explicit witnesses imposed on the constructed locus.

\begin{remark}[Bott periodicity and the framing obstruction]
\label{rem:v3-st-bott-obstruction}
The obstruction to extending the $\Sd$-framing beyond the
base $\Eone$ level is controlled by $\pi_d(BU)$:
\begin{center}
\renewcommand{\arraystretch}{1.2}
\small
\begin{tabular}{lllll}
 \toprule
 $d$ & Native $\En$ & $\pi_d(BU)$ & Obstruction & Example \\
 \midrule
 $1$ & $E_\infty$ & $0$ & trivial & elliptic curve \\
 $2$ & $\Etwo$ & $\bZ$ & braiding ($c_1$) & K3, abelian surface \\
 $3$ & $\Eone$ & $0$ & trivial & $K3 \times E$, conifold \\
 $4$ & $\Eone$ & $\bZ$ & shifted symplectic ($p_1$) & $K3 \times K3$ \\
 \bottomrule
\end{tabular}
\end{center}
The complex Bott group $\pi_d(BU)=KU^{-d}$ is $2$-periodic:
\[
 \pi_d(BU)=\begin{cases}\bZ,&d\ \mathrm{even},\\0,&d\ \mathrm{odd}.
 \end{cases}
\]
The $8$-periodicity belongs to the real and symplectic refinements
$\pi_d(BO)$ and $\pi_d(B\Sp)$.  In the Vol~III effective obstruction
tower the trivial dimensions occur at $d\bmod 8\in\{1,3,7\}$, with
a separate $\bZ/2$ symplectic refinement at $d\equiv5\pmod 8$.
\end{remark}


%%% ================================================================
%%% 3. THE FUNCTOR PHI
%%% ================================================================

\section{The CY-to-chiral functor}
\label{sec:functor}

\subsection{The two-stage factorisation and its four-step
realisation}

On the verified loci, the CY-to-chiral assignment factorises as
\[
 \Phi_d \;=\; \SpCh_{\Sigma_{d-1}, C} \circ
 \Phi^{\mathrm{FA}}_d
 \colon \CY_d^{\mathrm{wit}}\text{-}\mathrm{Cat} \longrightarrow
 \En\text{-}\mathrm{ChirAlg}.
\]
Stage~1 produces the canonical $E_d$-holomorphic factorisation
algebra $\Phi^{\mathrm{FA}}_d(\cC)$ on the CY$_d$ variety~$X$ after
the relevant formality and holomorphic-locality witnesses have been
fixed.  Kontsevich--Tamarkin $E_d$-formality acts on Hochschild
cochains with their degree-$-1$ Gerstenhaber bracket; after CY duality
to homology the shifted bracket has degree $1-d$. Stage~2 takes the
factorisation
homology of $\Phi^{\mathrm{FA}}_d(\cC)$ over a closed
$(d-1)$-cycle $\Sigma_{d-1} \subset X$ and restricts to a
reference curve $C \subset X$. A single CY$_d$ category thus
admits a family of $\Eone$-chiral shadows indexed by
$(\Sigma_{d-1}, C)$.

The four-step construction below is the computational incarnation
of the composite $\SpCh_{\Sigma_{d-1}, C} \circ
\Phi^{\mathrm{FA}}_d$ along a codimension-one cycle followed by
curve restriction; each step consumes a specific piece of the CY
data and produces a specific algebraic structure on the reference
curve~$C$.

\begin{construction}[The functor $\Phi$]
\label{con:v3-st-phi}
Let $\cC$ be a smooth proper CY$_d$ category equipped with the
negative-cyclic CY class, formality, holomorphic-twist, framing,
completion, and specialisation witnesses required in the dimension
under discussion.  Let $C$ be the reference curve.
\begin{enumerate}[label=\textbf{\arabic*.}]
 \item \textbf{Cyclic $\Ainf \to$ Lie conformal algebra.}
 The Gerstenhaber bracket on $\HH^\bullet(\cC)$ is a graded Lie
 bracket of degree~$-1$ (Proposition~\ref{prop:v3-st-gerstenhaber}).
 The CY pairing promotes this to a Lie conformal algebra
 $\fL_\cC$: the bracket becomes a $\lambda$-bracket
 $\{a_\lambda b\}$, and the pairing becomes the invariant
 bilinear form.  This step consumes the $\Ainf$-structure and
 the CY trace; it produces the current algebra of $\cC$.

 \item \textbf{Lie conformal algebra $\to$ factorization
 envelope.}
 The factorization envelope construction produces a factorization
 algebra $\Fact_X(\fL_\cC)$ on any smooth curve~$X$.  This step
 consumes the Lie conformal algebra and the curve; it produces a
 cosheaf on $\Ran(X)$ whose sections are the OPE data.  The
 universal enveloping vertex algebra $U_X^{\mathrm{fact}}(\fL_\cC)$
 is the chiral algebra.

 \item \textbf{$\Sd$-framing $\to$ $\En$ enhancement.}
 The framing witness on $\HH_\bullet(\cC)$
 (Proposition~\ref{prop:v3-st-sd-framing}) enhances the
 factorization algebra to an $\En$-chiral algebra.  At $d = 2$:
 the $\bS^2$-framing of Kontsevich--Vlassopoulos gives an
 $\Etwo$-chiral algebra.  At $d \geq 3$: the native $E_d$-structure
 restricts to $\Eone$ for CY quantum group purposes (the direct
 $E_d$-braiding is over-symmetric for braided chiral purposes; the
 ordered $\Eone$ direction is retained on~$C$).

 \item \textbf{Quantization.}
 The negative-cyclic CY class determines the quantization datum for
 the classical factorization algebra; its Hochschild trace is only the
 decategorified shadow.  The
 quantum chiral algebra $A_\cC = \Phi(\cC)$ is the output.  At
 $d = 2$, the Serre functor $\bS_\cC \simeq [2]$ kills the
 one-loop anomaly by parity, so the quantization preserves
 $\kappa_{\mathrm{ch}}$.  At $d = 3$, the Serre functor
 $\bS_\cC \simeq [3]$ no longer guarantees this; quantum
 corrections may shift $\kappa_{\mathrm{ch}}$.
\end{enumerate}
\end{construction}

\subsection{Theorem CY-A: the $d = 2$ case}

\begin{theorem}[CY-to-chiral functor at $d = 2$]
\label{thm:v3-st-cy-a2}
There exists a functor
$\Phi \colon \CY_2\text{-}\mathrm{Cat} \to
\Etwo\text{-}\mathrm{ChirAlg}$
from smooth proper CY categories of dimension~$2$ to
$\Etwo$-chiral algebras, such that:
\begin{enumerate}[label=(\roman*)]
 \item $\Phi(\cC)$ is a chiral algebra whose generating fields are
 in bijection with $\HH^{\bullet+1}(\cC)$.
 \item The bar complex $B(\Phi(\cC))$ is quasi-isomorphic to the
 cyclic bar complex $\CC_\bullet(\cC)$ as a factorization
 coalgebra.
 \item The modular characteristic is the chiral trace supertrace
 \[
  \kappa_{\mathrm{ch}}(\Phi(\cC))=\Xi(\cC).
 \]
 For $\cC=D^b(\Coh(X))$ with $h^{1,0}(X)=0$,
 \[
  \Xi(X)=\sum_q(-1)^q h^{0,q}(X)=\chi(\cO_X).
 \]
 The full Hochschild Euler characteristic
 $\sum_i(-1)^i\dim\HH_i(\cC)$ is a different invariant.
\end{enumerate}
\end{theorem}

\begin{proof}
The first two parts of Construction~\ref{con:v3-st-phi} are standard: the
Gerstenhaber bracket produces a Lie conformal algebra, and the
factorization envelope produces a chiral algebra.  The framing part at
$d = 2$ uses the Kontsevich--Vlassopoulos $\bS^2$-framing on
$\HH_\bullet(\cC)$, which gives an $\Etwo$-algebra structure on
the cyclic homology; this enhances the factorization algebra to an
$\Etwo$-chiral algebra.  The CY trace supplies quantization; at
$d = 2$, the Serre functor $\bS_\cC \simeq [2]$ forces the
holomorphic anomaly to vanish on the stated Hodge-filtered trace
lane.  Part~(i) follows from the construction; part~(ii) from the
identification of the bar differential with the cyclic boundary
operator; part~(iii) from the Hodge-filtered genus-one coefficient.
\end{proof}

\begin{remark}[Verification at $d = 2$]
\label{rem:v3-st-d2-verification}
For K3 surfaces: $\Xi(D^b(\Coh(S))) = \chi(\cO_S) = 2$,
and $\kappa_{\mathrm{ch}} = 2$ is verified by six independent
paths (direct Hodge, Mukai pairing, Dolbeault contracting
homotopy, bar computation, Borcherds lift, elliptic genus).  For
abelian surfaces: $\kappa_{\mathrm{ch}} = \chi(\cO_A) = 0$.
\end{remark}

\begin{conjecture}[Compact Hodge-trace lane at $d \geq 3$]
\label{conj:v3-st-cy-kappa-d3}
For $\cC=D^b(\Coh(X))$ on a constructed compact Hodge-trace lane with
$d \geq 3$,
\[
 \kappa_{\mathrm{ch}}(\Phi(\cC)) \;=\; \Xi(X)=\chi(\cO_X).
\]
This conjecture does not identify the compact Hodge supertrace with
the Heisenberg--Mukai specialisation or with the BKM weight.
\end{conjecture}

\begin{remark}[Evidence]
\label{rem:v3-st-d3-evidence}
At $d = 3$, compact total-space values such as
$\Xi(K3\times E)=0$ are distinct from the Heisenberg--Mukai value
$3$ and the BKM value $5$.  Non-compact toric local models require
their own normalization.  The obstruction to extending the $d = 2$
proof is that the $d = 3$ quantization step may introduce corrections:
Serre duality $\bS_\cC \simeq [3]$ no longer kills the one-loop
anomaly by parity.
\end{remark}

\subsection{The five objects in the CY context}

Under the functor $\Phi$, the five-object chain of Volume~I
becomes a chain of CY invariants.  The five objects are produced by
four distinct functors and should never be conflated.

\begin{remark}[Five objects from $\Phi(\cC)$]
\label{rem:v3-st-five-objects}
Let $A_\cC = \Phi(\cC)$ be the chiral algebra of a CY$_d$
category~$\cC$.
\begin{enumerate}[label=(\roman*)]
 \item $A_\cC$ (the chiral algebra): the factorization envelope
 of the Lie conformal algebra $\fL_\cC$.
 \item $B(A_\cC) = T^c(s^{-1}\bar{A}_\cC)$ (the bar coalgebra):
 a factorization coalgebra on $\Ran(X)$ with deconcatenation
 coproduct.  At genus $g \geq 1$, the fiberwise differential
 satisfies $d_{\mathrm{fib}}^2 = \kappa_{\mathrm{ch}}(A_\cC)
 \cdot \omega_g$.
 \item $A_\cC^i = H^*(B(A_\cC))$ (the dual coalgebra): its
 coalgebra structure encodes the BPS spectrum of~$\cC$.
 \item $A_\cC^! = ((A_\cC^i)^\vee)$ (the Koszul dual algebra):
 obtained by Verdier duality $D_{\Ran}$.  For CY$_2$ categories,
 $A_\cC^!$ is the chiral algebra of the mirror CY$_2$.
 \item $\Zder(A_\cC)$ (the derived chiral center):
 the chiral Hochschild cochain complex encoding the bulk sector.
\end{enumerate}
The bar complex is the $\Eone$ engine; the derived center is the
$\SCchtop/E_3$ output.  Bar-cobar
$\Omega(B(A_\cC)) \xrightarrow{\sim} A_\cC$ is inversion
(recovering the algebra), not Koszul duality; the derived center
is the categorical Hochschild construction, not the output of bar-cobar.
\end{remark}


%%% ================================================================
%%% 4. THE KAPPA-SPECTRUM
%%% ================================================================

\section{The $\kappa_\bullet$-spectrum}
\label{sec:kappa-spectrum}

\subsection{Four modular characteristics}

A CY manifold~$X$ admits multiple chiral algebraizations, each
with its own modular characteristic.  The four approved subscripts
in this volume are: $\kappa_{\mathrm{cat}}$ (categorical Hodge
supertrace), $\kappa_{\mathrm{ch}}$ and its explicit specialisations
such as $\kappa_{\mathrm{ch}}^{\mathrm{Heis}}$ (boundary chiral
trace lanes), $\kappa_{\mathrm{BKM}}$ (Borcherds--Kac--Moody), and
$\kappa_{\mathrm{fiber}}$ (lattice).
Bare $\kappa$ without subscript is forbidden.

\begin{definition}[The $\kappa_\bullet$-spectrum]
\label{def:v3-st-kappa-spectrum}
For a CY$_d$ manifold~$X$ (or more generally a CY$_d$
category~$\cC$), the \emph{$\kappa_\bullet$-spectrum} is the
ordered tuple
\[
 \Spec_{\kappa_\bullet}(X)
 = (\kappa_{\mathrm{cat}},\; \kappa_{\mathrm{ch}}^{\mathrm{Heis}},\;
 \kappa_{\mathrm{BKM}},\; \kappa_{\mathrm{fiber}})
\]
where:
\begin{enumerate}[label=(\roman*)]
 \item $\kappa_{\mathrm{cat}} = \chi(\cO_X)$
 (the holomorphic Euler characteristic) measures the total-space
 categorical Hochschild datum.  On products it is
 K\"unneth-multiplicative; fibre multipliers such as
 $\chi(\cO_{K3})=2$ must be named on the fibre lane, not imported into
 $\kappa_{\mathrm{cat}}(K3\times E)$.
 \item $\kappa_{\mathrm{ch}}^{\mathrm{Heis}}$ is the
 Heisenberg--Mukai boundary specialisation.  For $K3\times E$ it is
 $3=\dim_\bC(K3\times E)$.  The compact total-space chiral trace lane
 is instead $\kappa_{\mathrm{ch}}(A_{K3\times E})=\Xi(K3\times E)=0$.
 For $d = 2$ and $h^{1,0}=0$, Theorem~\ref{thm:v3-st-cy-a2} gives
 $\kappa_{\mathrm{ch}}=\chi(\cO_X)$.
 \item $\kappa_{\mathrm{BKM}}=c_N(0)/2$ measures the bulk
 Borcherds--Kac--Moody superalgebra in the Vol~III normalization.
 For $K3 \times E$, the primitive denominator is $\Delta_5$ and
 $\kappa_{\mathrm{BKM}}=5$; the squared unnormalised Igusa convention is
 $\Phi_{10}^{\mathrm{un}}=\Delta_5^2$.
 \item $\kappa_{\mathrm{fiber}} = \operatorname{rank}(\Lambda_X)$
 (the lattice rank) measures the frame data.  For $K3 \times E$:
 $\kappa_{\mathrm{fiber}} = 24$.
\end{enumerate}
\end{definition}

The four characteristics are generically distinct.  Their relations
encode the internal structure of the CY manifold.

\begin{proposition}[Spectrum relations for $K3 \times E$]
\label{prop:v3-st-k3e-spectrum}
For $X = K3 \times E$:
\[
 \Spec_{\kappa_\bullet}(K3 \times E)
 = (0,\; 3,\; 5,\; 24).
\]
Here the first entry is the total-space value
$\chi(\cO_{K3\times E})=0$, while the value $2$ belongs only to the
K3 fibre.  No additive identity equates the Borcherds weight with
the sum of the categorical and chiral entries.
The ratio
$\kappa_{\mathrm{BKM}}/\kappa_{\mathrm{ch}}^{\mathrm{Heis}} = 5/3$ reflects the
single-copy-to-second-quantization passage between the Heisenberg
specialisation and the Borcherds denominator.
\end{proposition}

\begin{proof}
$\kappa_{\mathrm{cat}} = \chi(\cO_{K3\times E})
= \chi(\cO_{K3})\chi(\cO_E)=2\cdot 0=0$ by K\"unneth multiplicativity
for the total space.
$\kappa_{\mathrm{ch}}^{\mathrm{Heis}} = \dim_\bC(K3 \times E) = 3$ (additivity
under Cartesian product, Route~B:
$\kappa_{\mathrm{ch}}^{\mathrm{Heis}}(K3) = 2$,
$\kappa_{\mathrm{ch}}^{\mathrm{Heis}}(E) = 1$;
verified by six independent paths).  This is the Heisenberg-level
additive invariant; the total-space Hodge-supertrace Route~A
$\Xi(X)=\chi(\cO_X)$ is K\"unneth-multiplicative and vanishes on
$K3 \times E$.
The Volume~III Beauville-route comparison separates these two
invariants: Route~A is the total-space Hodge supertrace, while
Route~B is the additive Heisenberg-level shadow.
$\kappa_{\mathrm{BKM}} = c_1(0)/2=\operatorname{wt}(\Delta_5)
=\frac{1}{2}\operatorname{wt}(\Phi_{10}^{\mathrm{un}})=5$ because
$\Phi_{10}^{\mathrm{un}} = \Delta_5^2$ has weight~$10$.
$\kappa_{\mathrm{fiber}} = \operatorname{rank}(H^*(K3, \bZ))
= 24$ (the second Betti number $b_2(K3) = 22$ plus
$b_0 + b_4 = 2$).  The numerical decomposition $5=2+3$ uses the K3
fibre value $2$ and the Heisenberg specialisation $3$; it is not an
identity among total-space $\kappa_\bullet$-invariants.  The fibre
rank relation $24 = 12 \cdot 2$ uses $\chi(\cO_{K3})=2$, not
$\chi(\cO_{K3\times E})$.
\end{proof}

\subsection{Koszul complementarity in the CY context}

Volume~I defines the Koszul conductor $K(A) =
\kappa_{\mathrm{ch}}(A) + \kappa_{\mathrm{ch}}(A^!)$: it is $0$
for KM/Heisenberg/lattice/free families, $13$ for Virasoro, and
family-dependent in general.  The CY context introduces a new
complementarity: mirror symmetry exchanges $\cC$ and $\cC^\vee$,
and $A_{\cC^\vee} = A_\cC^!$ when both are constructed.

\begin{conjecture}[CY Koszul complementarity]
\label{conj:v3-st-cy-koszul-complementarity}
For a CY$_2$ category $\cC$ with mirror $\cC^\vee$:
\[
 \kappa_{\mathrm{ch}}(\Phi(\cC)) + \kappa_{\mathrm{ch}}
 (\Phi(\cC^\vee)) = 0.
\]
The Koszul conductor vanishes for all mirror pairs at $d = 2$.
\end{conjecture}

\begin{remark}[Scope]
\label{rem:v3-st-cy-complement-evidence}
The quintic threefold is a CY$_3$ example and does not supply
evidence for the CY$_2$ complementarity statement.  Moreover
$\chi(\cO_{X_5})=0$ and $\chi_{\mathrm{top}}(X_5)=-200$, so
$\chi(X_5)/2$ is not the categorical trace invariant used here.
The CY$_3$ BCOV and mirror lane is a separate problem.
\end{remark}

\subsection{Shadow depth classification}

The CY-to-chiral functor assigns to each CY geometry a chiral
algebra whose shadow class (Volume~I) organizes the landscape.

\begin{proposition}[Shadow depth of CY chiral algebras]
\label{prop:v3-st-cy-shadow-depth}
Under $\Phi$, the shadow depth classification of Volume~I
(G/L/C/M) distributes over the CY landscape as follows:
\begin{center}
\renewcommand{\arraystretch}{1.3}
\small
\begin{tabular}{llll}
 \toprule
 CY geometry & Shadow class & $\kappa_{\mathrm{ch}}$ &
 Shadow depth \\
 \midrule
 $\bC^3$ (toric, abelian) & G & $1$ & $2$ \\
 Resolved conifold & G & $1$ & $2$ \\
 Local $\bP^2$ & M & $3/2$ & $\infty$ \\
 $K3 \times E$ & M & $3$ & $\infty$ \\
 \bottomrule
\end{tabular}
\end{center}
Classes~G and~M are populated.  Classes~L (shadow depth~$3$) and~C
(shadow depth~$4$) are targets of the toric programme: the toric
vertex bound $r_{\max} \leq N$ (where $N$ is the number of
web-diagram vertices) predicts that specific toric geometries
should realize these classes, but explicit CY examples at shadow
depth exactly~$3$ or~$4$ remain to be identified.
\end{proposition}


%%% ================================================================
%%% 5. K3 x E PROTOTYPE (PARTIAL -- SECTION HEAD)
%%% ================================================================

\section{The $K3 \times E$ prototype}
\label{sec:k3e}

The product of a K3 surface~$S$ and an elliptic curve~$E$ is the
canonical CY$_3$ testing ground.  It is simple enough to compute
(lattice $\Lambda^{3,2}$ of signature $(3,2)$, K3 elliptic genus
$\phi_{0,1}(\tau, z)$, Borcherds denominator $\Delta_5$) and
already exhibits every obstruction that separates the shadow tower
from the full automorphic form.

\subsection{The thirteen structural results}

The structural spine begins with six computations: modular
characteristic, moonshine multiplier, BKM weight, BPS factorisation,
mock modular decomposition, and shadow convergence.

\emph{K3-1} (Heisenberg-level modular characteristic, Route~B):
$\kappa_{\mathrm{ch}}^{\mathrm{Heis}}(K3 \times E) = 3 = \dim_\bC$,
verified by six independent paths (direct Hodge, Mukai pairing,
Dolbeault contracting homotopy, bar computation, Borcherds lift,
elliptic genus).  Additivity under Cartesian product:
$\kappa_{\mathrm{ch}}^{\mathrm{Heis}}(K3) = 2$,
$\kappa_{\mathrm{ch}}^{\mathrm{Heis}}(E) = 1$.  The Route~A
Hodge-supertrace $\Xi(X)=\chi(\cO_X)$ is K\"unneth-multiplicative and gives
$\Xi(K3 \times E) = 2 \cdot 0 = 0$; the two routes label
genuinely distinct invariants.  The Volume~III Beauville-route
comparison records the same split between total-space trace and
Heisenberg-level additive shadow.

\emph{K3-2} (moonshine multiplier): the factor~$2$ in
$Z_{K3} = 2\,\phi_{0,1}$ is
$\kappa_{\mathrm{cat}}(K3) = \chi(\cO_{K3}) = 2$.  The massive
Mathieu multiplicities satisfy
$A_n = \kappa_{\mathrm{cat}} \cdot \dim(\rho_n)$: the bar complex
provides the universal coefficient and does not detect~$M_{24}$
directly.

\emph{K3-3} (BKM weight):
$\kappa_{\mathrm{BKM}} = \frac{1}{2}\operatorname{wt}(\Phi_{10}^{\mathrm{un}})
= 5$.  The ratio
$\kappa_{\mathrm{BKM}}/\kappa_{\mathrm{ch}} = 5/3$ reflects
second quantization; the Hilbert--Chow exceptional divisor accounts
for $\chi(\operatorname{Hilb}^2) - \chi(\operatorname{Sym}^2) =
24 = \chi(K3)$.

\emph{K3-4} (BPS factorization):
$A_n = \kappa_{\mathrm{cat}}(K3) \cdot \dim(\rho_n)$, separating the
fibre homological part ($\kappa_{\mathrm{cat}}(K3) = 2$) from the
group-theoretic part ($\rho_n$, the Mathieu representation).

\emph{K3-5} (mock modular verification): the mock modular
decomposition
$Z_{K3} = (h - 24\mu)\,\vartheta_1^2/\eta^3$ is verified to
machine precision ($\sim 5.7 \times 10^{-16}$ relative error),
where $\mu$ is the Zwegers Appell--Lerch sum.

\emph{K3-6} (shadow convergence): the shadow generating function
$Z^{\mathrm{sh}}(t) = \kappa_{\mathrm{ch}}\bigl((t/2)/
\sin(t/2) - 1\bigr)$ converges with radius $R = 2\pi$, has entire
Borel transform, and exhibits no resurgence.

\subsection{The shadow--Siegel gap theorem}

\begin{theorem}[Shadow--Siegel gap]
\label{thm:v3-st-shadow-siegel-gap}
The shadow obstruction tower of $K3 \times E$ does not produce the
unnormalised Igusa cusp form $\Phi_{10}^{\mathrm{un}}$.  Four
structural obstructions
separate the two:
\begin{enumerate}[label=\textup{(O\arabic*)}]
 \item \emph{Categorical}: the shadow tower produces numbers
 (Taylor coefficients of $\Theta_{A_\cC}$); the denominator
 identity is a function (an automorphic form on~$\cH_2$).
 \item \emph{$\kappa_{\mathrm{ch}}^{\mathrm{Heis}}$--$\kappa_{\mathrm{BKM}}$
 mismatch}:
 \[
  \kappa_{\mathrm{ch}}^{\mathrm{Heis}}=3\neq 5=\kappa_{\mathrm{BKM}}.
 \]
 The compact total-space Hodge-supertrace value is separately
 $\kappa_{\mathrm{ch}}(A_{K3\times E})=0$.  The number $3$ is the
 Heisenberg--Mukai specialisation, not the total-space $\Phi$-trace.
 \item \emph{Second quantization}: the physical DT partition
 function is the DMVV symmetric product of single-copy data; the
 bar complex of a single algebra is the single copy.
 \item \emph{Schottky}: at $g \geq 4$, the Torelli map
 $\cM_g \to \cA_g$ is no longer dominant, and the shadow tower
 is imprisoned in tautological classes on $\Mbar_g$.
\end{enumerate}
\end{theorem}

Each obstruction points to a research programme.  The
$\kappa_{\mathrm{ch}}$--$\kappa_{\mathrm{BKM}}$ mismatch~(O2) is
the most telling: a single CY$_3$ admits multiple chiral
algebraizations, and the $\kappa_\bullet$-spectrum records which
face of the geometry each algebraization captures.


%%% ================================================================
%%% 6. TORIC CY3 AND CoHA (SECTION HEAD)
%%% ================================================================

\section{Toric CY$_3$ and the cohomological Hall algebra}
\label{sec:toric}

A toric CY$_3$ $X_\Sigma$ is determined by a fan~$\Sigma$ in
$\bZ^3$ satisfying the CY condition (all generators coplanar).
The toric diagram fixes a quiver with potential $(Q_X, W_X)$, and
the critical CoHA
$\cH(Q_X, W_X) = \bigoplus_\mathbf{d}
H^{\mathrm{BM}}_*(\mathrm{Crit}(W_\mathbf{d}),
\phi_{W_\mathbf{d}})$ provides the $\Eone$-sector of the CY
quantum group.

\subsection{$\bC^3$ and the affine Yangian}

The simplest toric CY$_3$.  Toric diagram: a single triangle.
Quiver: the Jordan quiver (one vertex, one loop~$\circlearrowleft$).
Cubic potential $W = \mathrm{tr}(XYZ - XZY)$.

\begin{theorem}[Schiffmann--Vasserot]
\label{thm:v3-st-sv}
The critical CoHA of $\bC^3$ is isomorphic to the positive half of
the affine Yangian of $\fgl_1$:
\[
 \CoHA(\bC^3) \;\simeq\; Y^+(\widehat{\fgl}_1).
\]
The full $\cW_{1+\infty}$ vertex realisation appears only after
taking the appropriate double/evaluation realisation.
\end{theorem}

The chain $\bC^3 \to Y^+(\widehat{\fgl}_1) \to
\Rep^{\Etwo}(Y(\widehat{\fgl}_1))$ is verified end-to-end: the
toric CY$_3$ produces the CoHA, which is the positive half of the
Yangian and hence the $\Eone$-chiral quantum group.  The
$\cW_{1+\infty}$ vertex realisation enters only after the
double/evaluation step.  The Drinfeld center recovers the full
$\Etwo$-braiding
(Theorem~\ref{thm:v3-st-drinfeld-center} below).

\subsection{The CoHA as $\Eone$-primitive}

The CoHA is an associative ($\Eone$) algebra: the Hall
multiplication is ordered (short exact sequences
$0 \to V' \to V \to V'' \to 0$ have sub before quotient).  In the
present framework:

\begin{itemize}
 \item The CoHA is the positive half: the $\Eone$-chiral sector
 (the ordered part).
 \item The full quantum vertex realisation is obtained only after the
 double/center construction, where the $\Etwo$-braiding lives.
 \item The braiding (the passage from $\Eone$ to $\Etwo$) is
 the quantum group $R$-matrix of the affine super Yangian.
 \item The ordered bar $B^{\mathrm{ord}}(A_\cC)$ preserves the
 $R$-matrix; the symmetric bar $B^\Sigma(A_\cC)$ kills the Hopf
 structure via averaging: in abelian and scalar families
 $\mathrm{av}(r(z)) = \kappa_{\mathrm{ch}}$, while for
 non-abelian affine Kac--Moody
 $\mathrm{av}(r(z))+\dim(\fg)/2 = \kappa_{\mathrm{ch}}$.
\end{itemize}

For local toric charts the hCS BV complex and the Hall shuffle algebra
glue at finite positive-half Rees level: after truncating the Rees
filtration, the positive-half Hall multiplication and the hCS extension
correspondence define the same ordered $\Eone$ operation.  Vol~III now
constructs this passage on the resolved conifold by gluing the hCS
$E_3$ factorisation algebra, adding the anomaly-free quantum lift, and
writing the finite hCS--Rees-Hall chart maps.  The compact Hall,
quasi-NCCR, and recognised Hall-double comparisons are separate global
identifications; finite radical-quotient doubles now exist from compact
windows, while the Borcherds recognition/completion remains a
finite-defect and inverse-limit condition.

\begin{remark}[Shadow = $\GW(\bC^3)$]
\label{rem:v3-st-shadow-gw}
At $\kappa_{\mathrm{ch}} = \Psi$ (the deformation parameter of
$\cW_{1+\infty}$ at $c = 1$), the shadow obstruction tower recovers
the genus-$1$ constant-map channel for $\bC^3$. The all-genus virtual
identification requires the uniform-weight and $\Phi_3$ hypotheses,
and the motivic identification remains conjectural. On the DT side,
the MacMahon function
$M(q) = \prod_{n \geq 1}(1 - q^n)^{-n}$ encodes the partition
function as the second-quantized bar/DT character in the non-compact
degree-zero normalization.
\end{remark}

\subsection{The Drinfeld center}

\begin{theorem}[Drinfeld center for $\bC^3$]
\label{thm:v3-st-drinfeld-center}
The Drinfeld center of the $\Eone$-monoidal representation
category of $Y^+(\widehat{\fgl}_1)$ recovers the full
$\Etwo$-braided monoidal structure:
\[
 \cZ(\Rep^{\Eone}(Y^+(\widehat{\fgl}_1)))
 \;\simeq\;
 \Rep^{\Etwo}(Y(\widehat{\fgl}_1))
 \;\simeq\;
 \Rep^{\Etwo}(\cW_{1+\infty}).
\]
Five invariants match at six levels: the ordinary center has
dimension~$1$, while the derived center has dimension~$3$
(capturing $\Ext^{1,2}$ invisible to the commutant).
\end{theorem}

The Drinfeld center is the categorified averaging map
$\mathrm{av} \colon \fg^{\Eone} \to \fg^{\mathrm{mod}}$
of Volume~I: it projects the ordered $R$-matrix data to the scalar
$\kappa_{\mathrm{ch}}$ at the level of the $\kappa_\bullet$-spectrum,
while recovering non-trivial braiding at the categorical level.


%%% ================================================================
%%% 7. E_1 STABILIZATION (SECTION HEAD)
%%% ================================================================

\section{The $\Eone$ stabilization theorem}
\label{sec:e1-stabilization}

\begin{theorem}[$\Eone$ stabilization for toric local models]
\label{thm:v3-st-e1-stabilization}
For toric $GL(d)$-invariant local CY models with $d \geq 3$, the
Stage-2 CY chiral algebra $A_\cC=\Phi(\cC)$ is natively $\Eone$ on
the reference curve.  In this local invariant lane the selected CY
polyvector sector has trivial Schouten--Nijenhuis bracket, and the
CoHA extension correspondence imposes a preferred ordering (sub
before quotient).  For general CY categories the full
Schouten--Nijenhuis bracket remains nonzero; the $\Eone$ level is
forced by Stage~2 and degree.  The braided monoidal structure is
recovered through the Drinfeld center.
\end{theorem}

\begin{proof}
The Schouten--Nijenhuis bracket on $\mathrm{PV}(\bC^d)$ is nonzero in
general; vector fields already bracket.  The toric local model uses
the $GL(d)$-invariant CY sector generated by the volume class and the
scalar torus-weight data entering the Hall construction.  On this
selected sector the Schouten output has no invariant target of the
required CY weight, hence the induced bracket is zero.  The quantum
operation then comes from the CoHA convolution correspondence, which
is associative ($\Eone$) and ordered by subobject before quotient.
\end{proof}

\begin{remark}[Arithmetic of the fundamental group]
\label{rem:v3-st-pi1-arithmetic}
The key arithmetic for ordered configurations is
\[
\pi_1(\Conf^{\mathrm{ord}}_2(\bR^2))=\bZ,\qquad
\pi_1(\Conf^{\mathrm{ord}}_2(\bR^d))=0\quad(d\geq3).
\]
For the unordered quotient,
\[
\pi_1(\Conf^{\mathrm{unord}}_2(\bR^d))=\bZ/2\quad(d\geq3).
\]
Thus the native $E_d$ operation at $d\geq3$ is over-symmetric for
braided chiral purposes; the construction keeps the ordered $\Eone$
direction on $C$ and recovers braiding through the Drinfeld center.
\end{remark}


%%% ================================================================
%%% Three consolidated structural theorems: native $\En$ output,
%%%   CY-dimension stratification, Borcherds weight
%%% ================================================================

\section{$\Phi$ scope, CY-dimension stratification, Borcherds weight}
\label{sec:phi-scope-stratification-trace}

Three consolidated structural theorems pin the CY-to-chiral
assignment to its precise scope: the native $\En$-output level on the
reference curve, the dimension-stratified trace spectrum, and the
Vol~III Borcherds weight normalisation
$\kappa_{\mathrm{BKM}} = c_N(0)/2$.

\begin{theorem}[Native $\En$-output scope of the two-stage $\Phi_d$]
\label{thm:phi-native-en-scope}
Let
\[
 \Phi_d^{(\Sigma_{d-1},C)}
 =\SpCh_{\Sigma_{d-1},C}\circ\Phi^{\mathrm{FA}}_d
\]
be evaluated on a constructed locus with the required formality,
holomorphic-twist, framing, completion, and specialisation witnesses.
Its native output level is
\[
 n \;=\; \begin{cases} \infty & d = 1, \\ 2 & d = 2, \\
1 & \text{constructed } d \ge 3. \end{cases}
\]
At $d = 2$ the native $\bS^2$-framing on $\HH_\bullet(\cC)$
supplies a genuine $\Etwo$-chiral structure
(Kontsevich--Vlassopoulos; Theorem~\textup{CY-A$_2$}).
At $d = 3$ this is a witnessed object-level assignment for fixed
$(\Sigma_2,C)$; arbitrary functoriality on all CY morphisms is not
asserted.  For $d\geq4$ the display names a construction problem, not
a theorem.  Braiding is recovered categorically via the Drinfeld
centre, closing the $\En$ operadic circle on
$Z^{\mathrm{der}}_{\mathrm{ch}}$ rather than on a bare output.
\end{theorem}

\begin{theorem}[Trace stratification by
 CY-dimension]\label{thm:phi-kappa-stratification}
For a compact Calabi--Yau manifold $X$ of dimension $d$ and a
reference curve $C \subset X$ with specialisation cycle
$\Sigma_{d-1} \subset X$,
\begin{equation}\label{eq:phi-kappa-hodge-supertrace}
\kappa_{\mathrm{ch}}^{\mathrm{Hodge}}\bigl(\Phi_d(D^b\Coh(X))\bigr)
\;=\; \sum_{q=0}^{d} (-1)^q h^{0,q}(X),
\end{equation}
on the compact Hodge-trace lane. The stratification across the canonical
families reads
\begin{center}\small
\renewcommand{\arraystretch}{1.2}
\begin{tabular}{lcl}
\toprule
family (dim $d$) & $\kappa_{\mathrm{ch}}^{\mathrm{Hodge}}$ & notes \\
\midrule
elliptic curve $E$ ($d=1$) & $0$ & $h^{0,0}-h^{0,1}=1-1$ \\
$K3$ ($d=2$) & $2$ & $\chi(\cO_{K3}) = 2$ \\
abelian/bielliptic surface ($d=2$) & $0$ & $\chi(\cO) = 0$ \\
quintic $\subset \bP^4$ ($d=3$) & $0$ & $\chi(\cO_X) = 0$ \\
$K3 \times E$ ($d=3$) & $0$ & $\chi(\cO) = 0$ \\
triple elliptic $E^3$ ($d=3$) & $0$ & $\chi(\cO) = 0$ \\
local $\bP^2 = \mathrm{Tot}(K_{\bP^2})$ & $3/2$
  & $\chi(S)/2 = 3/2$ (shadow) \\
$\mathrm{CY}_4$ sextic $\subset \bP^5$ & $2$
  & Hodge supertrace \\
generic $\mathrm{CY}_5$ & $0$ & Hodge supertrace vanishes \\
\bottomrule
\end{tabular}
\end{center}
The Heisenberg--Mukai, BKM, and fibre lanes are not entries of this
table.  The conifold is \emph{not} a local surface: its geometry is
$\mathrm{Tot}(\cO(-1) \oplus \cO(-1) \to \bP^1)$, not
$\mathrm{Tot}(K_S \to S)$, and the identity
$\kappa_{\mathrm{ch}}^{\mathrm{Hodge}} = \chi(S)/2$ does not apply. The McKay
orbifold $[\bC^3/\bZ_3]$ gives
$\kappa_{\mathrm{ch}}^{\mathrm{Hodge}} = 1$ by direct calculation.
\end{theorem}

\begin{proof}
The identification
$\kappa_{\mathrm{ch}}^{\mathrm{Hodge}} = \sum(-1)^q h^{0,q}$ is
the Hochschild--Kostant--Rosenberg degeneration composed with the
Mukai pairing; the degree-$d$ cyclic trace $\Tr\colon \HC^-_d(\cC) \to k$
restricts to the top-dimensional Hodge piece, and Serre duality
pairs it with the full Hodge decomposition. The family entries are
direct computation; ``local~$\bP^2$'' comes from the shadow
$\kappa_{\mathrm{ch}}^{\mathrm{Hodge}} = \chi(S)/2 = 3/2$ because
$\chi(\bP^2) = 3$. The conifold exclusion is structural: its
total space is not the total space of a canonical line bundle over
a surface, so the local-surface formula is inapplicable.
\end{proof}

\begin{theorem}[Borcherds weight normalisation]
\label{thm:phi-universal-trace}
On the Vol~III Borcherds-denominator locus
$N\in\{1,2,3,4,6\}$,
\[
 \kappa_{\mathrm{BKM}}(\Phi_N)=\frac{c_N(0)}{2}.
\]
Thus $c_N(0)=(10,8,6,4,2)$ gives
\[
 \kappa_{\mathrm{BKM}}=(5,4,3,2,1).
\]
This is a Borcherds-weight identity.  The compact Hodge supertrace is
separate, and the ghost-charge normalisation belongs to the
Monster--Leech lane.
\end{theorem}

\begin{remark}[$\kappa$ subscript conventions (all subscripts explicit)]
\label{rem:phi-kappa-subscripts}
Every $\kappa$ in this section carries an approved Vol~III
subscript from the closed set $\{\mathrm{ch}, \mathrm{cat},
\mathrm{BKM}, \mathrm{fiber}\}$. Bare $\kappa$ is
forbidden in Vol~III. The identifications
$\kappa_{\mathrm{ch}} \to \chi(\cO)$ (Hodge supertrace),
$\kappa_{\mathrm{ch}}^{\mathrm{Heis}} \to \dim_\bC X$ on the
Heisenberg--Mukai specialisation, and
$\kappa_{\mathrm{BKM}} \to c_N(0)/2$ are the precise uses here.
\end{remark}

\begin{remark}[Dimension range of $\Phi_d$]
\label{rem:phi-dimensional-range}
Stage~1 $\Phi^{\mathrm{FA}}_d$ is unconditional for $d \in \{1, 2\}$
and is constructed at $d=3$ on the witnessed filtered locus by
Kontsevich--Tamarkin $E_d$-formality applied to Hochschild cochains,
Costello--Gwilliam--Li locality on~$X$, and the required framing and
completion witnesses. Stage~2
$\SpCh_{\Sigma_{d-1}, C}$ is factorisation homology followed by
curve restriction along a witnessed admissible specialisation datum
$(\Sigma_{d-1},C)$.
$\Phi_d$ at $d = 1$: trivial ($\Ainf = \mathrm{comm}$).
$d = 2$: proved (Theorem~\textup{CY-A$_2$}), $\Etwo$-chiral
output. $d = 3$: witnessed object-level output for fixed
$(\Sigma_2,C)$, $\Eone$-chiral on the reference curve; the family of
constructed specialisations
$\{\SpCh_{\Sigma_2, C}(\Phi^{\mathrm{FA}}_3(\cC))\}_{\Sigma_2}$
over choices of $K3$-, $T^4$-, or Enriques-fibres produces
the Borcherds--$\Delta_5$, Igusa--$\chi_{10}$, and
Gritsenko--Cl\'ery paramodular shadows respectively.
$d = 4$: Stage~1 is obstructed by $\pi_4(BU) = \bZ$; the
$p_1$-twisted double current algebra is the surviving structure.
For $K3\times K3$ one has
$\int_{K3\times K3}p_1(T_X)^2=4608$, so this product lies on the
$p_1$-twisted stratum rather than on a Pontryagin-vanishing locus.
The untwisted $E_4$ lift requires separate vanishing data. For $d = 5, 6$ no general
Stage~1 theorem is asserted here; the expected Hodge-supertrace value
$\kappa_{\mathrm{ch}}=0$ for a generic smooth CY requires a separate
existence theorem for the relevant higher factorisation algebra.
Thus the unconditional range is $d \in \{1, 2\}$, the constructed
$d=3$ range is witnessed object-level, and the $d=4$
untwisted range is the Pontryagin-vanishing family; $p_1$-twisted
examples such as $K3\times K3$ require the twisted datum.
\end{remark}


%%% ================================================================
%%% REFERENCES
%%% ================================================================

\section*{References}

\begin{enumerate}[label={[\arabic*]}]
\item
 R.~Lorgat,
 \emph{Modular Koszul Duality},
 Volume~I of the Modular Koszul Duality Programme, 2026.

\item
 R.~Lorgat,
 \emph{$\Ainf$ Chiral Algebras and 3D Holomorphic-Topological QFT},
 Volume~II of the Modular Koszul Duality Programme, 2026.

\item
 R.~Lorgat,
 \emph{Calabi--Yau Categories, Quantum Groups, and BPS Algebras},
 Volume~III of the Modular Koszul Duality Programme, 2026.

\item
 K.~Costello,
 \emph{Notes on supersymmetric and holomorphic field theories in
 dimensions $2$ and $4$},
 Pure Appl.\ Math.\ Q.\ \textbf{9} (2013), 73--165.

\item
 K.~Costello and O.~Gwilliam,
 \emph{Factorization Algebras in Quantum Field Theory},
 Cambridge University Press, 2017, 2021.

\item
 M.~Kontsevich and Y.~Vlassopoulos,
 \emph{Pre-Calabi--Yau algebras and topological quantum field
 theories}, in preparation.

\item
 O.~Schiffmann and E.~Vasserot,
 \emph{Cherednik algebras, $\cW$-algebras and the equivariant
 cohomology of the moduli space of instantons on $\mathbb{A}^2$},
 Publ.\ Math.\ IHES \textbf{118} (2013), 213--342.

\item
 M.~Rapcak, Y.~Soibelman, Y.~Yang, and G.~Zhao,
 \emph{Cohomological Hall algebras, vertex algebras, and
 instantons},
 Comm.\ Math.\ Phys.\ \textbf{376} (2020), 1803--1873.

\item
 V.~Gritsenko, V.~Nikulin,
 \emph{Automorphic forms and Lorentzian Kac--Moody algebras},
 Part~II, Internat.\ J.\ Math.\ \textbf{9} (1998), 201--275.

\item
 R.~Dijkgraaf, E.~Verlinde, H.~Verlinde,
 \emph{Counting dyons in $N=4$ string theory},
 Nucl.\ Phys.\ B \textbf{484} (1997), 543--561.

\item
 D.~Maulik, N.~Nekrasov, A.~Okounkov, R.~Pandharipande,
 \emph{Gromov--Witten theory and Donaldson--Thomas theory, I, II},
 Compos.\ Math.\ \textbf{142} (2006), 1263--1285 and 1286--1304.
\end{enumerate}

\end{document}
